\chapter{Anatomy of a Makefile}\label{ch:mk-anatomy}

A Makefile is a list of rules. Everything else --- variables, functions,
pattern matching, conditionals --- exists to stop you from writing the same rule
forty times.

\section{The Rule}\label{sec:mk-rule}

\begin{center}
	\ttfamily\large
	\{target\}: [\{prerequisites\}]\\[2pt]
	\textrightarrow\{recipe\}
\end{center}

\begin{definition}[Target, prerequisite, recipe]
	The \vocab{target} is the file to be produced. The \vocab{prerequisites} are
	the files it is produced from. The \vocab{recipe} is the shell commands that
	produce it, each on its own line, each indented with a tab.
\end{definition}

\begin{makecode}
app: main.o parser.o util.o
    gcc -o app main.o parser.o util.o

main.o: main.c parser.h
    gcc -Wall -c main.c

parser.o: parser.c parser.h
    gcc -Wall -c parser.c

util.o: util.c
    gcc -Wall -c util.c
\end{makecode}

That is the graph of \autoref{sec:mk-graph} written out. Running \cmd{make app}
makes make read every rule, build the sub-graph rooted at \texttt{app}, and run
only the recipes whose targets are out of date.

\begin{shell}
$ make app
gcc -Wall -c main.c
gcc -Wall -c parser.c
gcc -Wall -c util.c
gcc -o app main.o parser.o util.o
$ make app                      # nothing changed
make: 'app' is up to date.
$ touch parser.c
$ make app                      # only what the graph reaches
gcc -Wall -c parser.c
gcc -o app main.o parser.o util.o
\end{shell}

\begin{note}
	Rules may appear in any order --- make reads the whole file before building
	anything, so a rule can name a prerequisite defined fifty lines later. The
	single exception is the \emph{first} rule in the file, which is special
	(\autoref{sec:mk-goal}).
\end{note}

\begin{note}[A target need not be a file]
	Nothing checks that \texttt{app} is really produced. A rule whose target is
	never created just runs its recipe every time, which is how \texttt{clean} and
	\texttt{test} work --- with a caveat that \autoref{sec:mk-phonyt} spends its
	whole length on.
\end{note}

\section{The Tab Character}\label{sec:mk-tab}

Every line of a recipe must begin with a literal tab. Not four spaces, not
eight: a tab.

\begin{center}
	\ttfamily
	\begin{tabular}{@{}l@{}}
		app: main.o                                                        \\
		{\color{WildStrawberry}\ensuremath{\hookrightarrow}}\quad gcc -o app main.o \\
	\end{tabular}
\end{center}

\begin{note}
	The arrow marks the one \texttt{U+0009} byte that must begin that line. This
	is make's most notorious wart --- its author called it a mistake within a few
	years, and by then it was too widely used to change.
\end{note}

Spaces instead of a tab produce make's least helpful error message:

\begin{shell}
$ make
Makefile:2: *** missing separator.  Stop.
\end{shell}

\begin{keys}[0.30]
	\cmd{cat -A Makefile}       & Show tabs as \texttt{\textasciicircum{}I} and line ends as \texttt{\$} \\
	\cmd{grep -P '\textbackslash t' Makefile} & Which lines actually contain a tab      \\
	\cmd{sed -n 'l' Makefile}   & The same, in POSIX                                    \\
	\key{<C-v>} \key{<Tab>}     & In Vim's Insert mode: insert a literal tab even where \cmd{'expandtab'} is on \\
\end{keys}

\begin{note}[Editors that expand tabs]
	Any editor configured to turn tabs into spaces will silently corrupt a
	Makefile. In Neovim, add a filetype rule (\autoref{sec:nvim-keymaps}):
\begin{luacode}
vim.api.nvim_create_autocmd("FileType", {
  pattern = "make",
  callback = function() vim.opt_local.expandtab = false end,
})
\end{luacode}
	LazyVim already does this, and \cmd{:set list} makes the tabs visible so you
	can check.
\end{note}

\begin{remark}
	GNU make since 3.82 can use another character instead. Setting the special
	variable \cmd{.RECIPEPREFIX} to \texttt{>} lets recipes be indented with
	\texttt{>} rather than a tab. It works, and almost nobody uses it: a Makefile
	that does is one more thing the next reader has to notice. Use the tab.
\end{remark}

\section{The Default Goal}\label{sec:mk-goal}

Running \cmd{make} with no arguments builds the \vocab{default goal}: the target
of the first rule in the file.

\begin{makecode}
all: app docs          # first rule -- so `make' means `make all'
app: main.o
    gcc -o app main.o
docs: manual.pdf
\end{makecode}

\begin{moral}
	Make the first rule \texttt{all}, always, and make it depend on everything a
	plain \cmd{make} should produce. Otherwise the default goal is whatever rule
	happens to be at the top, and adding a rule in the wrong place quietly changes
	what \cmd{make} does.
\end{moral}

\begin{note}
	\cmd{.DEFAULT\_GOAL := app} states it explicitly and is immune to reordering.
	It is worth using in a Makefile that other people will edit.
\end{note}

\section{How Recipes Are Run}\label{sec:mk-recipe}

\begin{moral}
	Each line of a recipe runs in its \emph{own} shell. Directory changes and
	variables do not carry to the next line.
\end{moral}

\begin{makecode}
broken:
    cd build          # this shell ends here...
    ls                # ...so this one is back where it started

works:
    cd build && ls    # one shell, one line

also-works:
    cd build; \
    ls                # a backslash joins them into one line
\end{makecode}

\begin{keys}[0.20]
	\texttt{@}   & Do not echo this line before running it            \\
	\texttt{-}   & Ignore a non-zero exit status; carry on            \\
	\texttt{+}   & Run even under \cmd{make -n} (\autoref{sec:mk-flags}) \\
\end{keys}

\begin{makecode}
greet:
    @echo "building..."    # prints the message, not the echo command
clean:
    -rm -f *.o             # do not fail if there is nothing to remove
\end{makecode}

\begin{note}
	Make aborts as soon as a recipe line exits non-zero, which is the behaviour you
	want --- a build that carries on after a compile error produces a binary built
	from stale objects. The \texttt{-} prefix switches it off for one line, and
	\cmd{rm} in a \texttt{clean} rule is the main legitimate use. \cmd{rm -f}
	is usually better, since it succeeds on its own.
\end{note}

\begin{remark}
	The shell is \file{/bin/sh}, not your login shell, so the Bash features of
	\autoref{ch:sh-script} --- \texttt{[[ ]]}, arrays, process substitution --- are
	unavailable unless you ask for them with \cmd{SHELL := /bin/bash}. A recipe
	that works when you paste it into your terminal and fails under make is almost
	always this.
\end{remark}

\section{A First Makefile}\label{sec:mk-first}

Everything so far, in one file, for a three-source C program:

\begin{makecode}
# ---- what to build -------------------------------------------------------
all: app

app: main.o parser.o util.o
    gcc -o app main.o parser.o util.o

main.o: main.c parser.h
    gcc -Wall -Wextra -g -c main.c

parser.o: parser.c parser.h
    gcc -Wall -Wextra -g -c parser.c

util.o: util.c util.h
    gcc -Wall -Wextra -g -c util.c

# ---- housekeeping --------------------------------------------------------
clean:
    rm -f app *.o

.PHONY: all clean
\end{makecode}

This is correct, and it is also the last time in Part~\ref{part:make} you will
see a Makefile this repetitive. Three things about it are unsatisfying:

\begin{enumerate}[(i)]
	\item The compiler and its flags are written out four times, so changing
	      \cmd{-O2} means four edits and one you will miss.
	\item The three \texttt{.o} rules are the same rule with different names.
	\item The header dependencies are maintained by hand, so the first time you add
	      an \cmd{\#include} and forget to record it, make will confidently not
	      rebuild.
\end{enumerate}

\autoref{ch:mk-vars} fixes the first, \autoref{ch:mk-patterns} the second and
third. The result is about fifteen lines and works for any number of sources.

\begin{tryit}
	Create the three source files (any contents will do), type this Makefile ---
	with real tabs --- and run \cmd{make}, then \cmd{make} again, then
	\cmd{touch parser.h} and \cmd{make}. Watch which recipes run each time, and
	check them against the graph on \autopageref{sec:mk-graph}.
\end{tryit}
